
La robustez de un sistema crítico como \textit{NodeAlert IoT} depende de su capacidad para gestionar múltiples eventos simultáneos con determinismo. Para lograr esto, se ha implementado el sistema operativo de tiempo real \textbf{FreeRTOS} sobre el microcontrolador ESP32, permitiendo una arquitectura de software basada en tareas concurrentes en lugar de un bucle secuencial tradicional \citep{valvano2014volume1}.

\subsection{Gestión de Tareas y Prioridades}
El diseño se basa en la división de funciones en seis tareas independientes, cada una con una prioridad y un tamaño de pila asignados según su relevancia para la seguridad del sistema \citep{valvano2014volume2}.

\begin{itemize}
    \item \textbf{taskKY026 (Prioridad 4, stack 2048 palabras):} Monitorea el sensor de llama KY-026 mediante interrupción por flanco de subida en GPIO5. Al detectar flama, notifica a \texttt{taskAutomation} mediante una cola FreeRTOS con timestamp. Es la tarea de mayor prioridad por la naturaleza crítica de un incendio.
    \item \textbf{taskAutomation (Prioridad 4, stack 2048 palabras):} Evalúa el estado global del sistema combinando las lecturas de los tres sensores. Aplica la máquina de estados con histéresis (ALERT\_NONE, WARNING, CRITICAL, EMERGENCY) y decide las acciones del actuador. Comparte prioridad máxima con taskKY026 para evitar cuellos de botella en la respuesta.
    \item \textbf{taskMQ9 (Prioridad 3, stack 2048 palabras):} Lee el sensor de gas MQ-9 cada 1 segundo mediante el ADC. Controla el calentador del sensor (GPIO18) con un tiempo de precalentamiento de 60 segundos. Publica el valor de gas en una cola para taskAutomation.
    \item \textbf{taskDHT22 (Prioridad 2, stack 3072 palabras):} Lee temperatura y humedad del DHT22 cada 2 segundos usando la biblioteca Adafruit DHT. Publica en un tópico MQTT separado y envía una copia a taskAutomation.
    \item \textbf{taskBuzzer (Prioridad 2, stack 1024 palabras):} Gestiona el buzzer en GPIO2 reproduciendo patrones PWM según la severidad:
    \begin{itemize}
        \item WARNING: 1 pitido cada 2 segundos.
        \item CRITICAL: 2 pitidos r\'apidos cada segundo.
        \item EMERGENCY: Tono continuo.
    \end{itemize}
    \item \textbf{taskMqttLoop (Prioridad 1, stack 4096 palabras):} Ejecuta el bucle de mantenimiento de PubSubClient (\texttt{mqttClient.loop()}) para mantener la conexión WiFi/MQTT, procesar mensajes entrantes y publicar telemetría. Es la tarea de menor prioridad porque su latencia no afecta la seguridad inmediata.
\end{itemize}

\subsection{Sincronización y Comunicación entre Tareas}
Para evitar condiciones de carrera y garantizar la integridad de los datos, se utilizan primitivas de sincronización de FreeRTOS. Según \citet{mazidi2014}, el intercambio de información entre procesos en arquitecturas de 32 bits debe ser atómico.
\begin{itemize}
    \item \textbf{Colas (Queues):} Se emplean dos colas: una para que los sensores envíen lecturas a la tarea de automatización, y otra para que la automatización envíe órdenes al gestor del buzzer.
    \item \textbf{Sem\'aforo Binario:} Se utiliza para notificar a la tarea KY026 desde la ISR (\textit{Interrupt Service Routine}) del GPIO5, siguiendo el patrón \texttt{xSemaphoreGiveFromISR} recomendado por FreeRTOS para despertar tareas desde interrupciones.
    \item \textbf{Variables Compartidas con Protección:} El estado global del sistema se almacena en una variable protegida por semáforo mutex, accesible desde taskAutomation y taskMqttLoop.
\end{itemize}

\subsection{Optimización}

En el diseño de sistemas IoT de misión crítica, la eficiencia en la comunicación y el procesamiento es fundamental para la sostenibilidad del sistema \citep{boylestad2009}.

\subsection{Optimización mediante FreeRTOS}
El ESP32 opera bajo FreeRTOS, que gestiona la CPU de forma eficiente bloqueando las tareas cuando están a la espera de eventos o temporizadores. Cuando todas las tareas están bloqueadas, FreeRTOS ejecuta la \texttt{tarea Idle}, que pone el núcleo en estado de bajo consumo. Esto ocurre naturalmente porque los sensores tienen periodos de muestreo largos (2 s para DHT22, 1 s para MQ-9, 200 ms para KY-026), y el CPU pasa la mayor parte del tiempo inactivo \citep{valvano2014volume1}.

\subsection{Eficiencia en la Comunicación}
La elección del protocolo MQTT contribuye directamente a la eficiencia del sistema. Al ser un protocolo orientado a eventos y de cabecera mínima, reduce el tiempo de activación de la radio WiFi (\textit{TX/RX time}), que es el componente de mayor consumo en el nodo \citep{mazidi2014}.